% Extensión estimada: unas 2.5 páginas, cuatro párrafos y una tabla (ficha 4.5).
% Fuentes: Tablas tab:marco_sota_clasificacion, tab:marco_sota_vlm y tab:marco_sota_vlm_gen del Capítulo 2 y las
% cifras de las secciones 4.1 a 4.4. Reúne las comparaciones con la literatura que antes estaban repartidas en cada
% sección (revisión del autor, 24/10/2026).
\section{Comparación con el estado del arte}
\label{sec:res_estado_arte}

Los trabajos del Capítulo~\ref{cp:marco} no permiten una comparación directa con estos resultados, porque difieren en
los datos, en la forma de dividirlos y en la regla de decisión. Por eso, esta sección no ordena los modelos, sino que
ubica el orden de magnitud de los resultados y examina qué parte de las diferencias se explica por el protocolo. La
Tabla~\ref{tab:res_estado_arte} reúne las cifras comparables en tres bloques: la clasificación en siete clases, la
detección de melanoma y los \glspl{vlm} sin ajuste. Cada fila es un trabajo, con sus datos y su protocolo, la métrica
que se compara, el valor que informa y el del modelo equivalente de este trabajo.

\begin{table}[H]
    \centering
    \caption{Resultados de la literatura (Capítulo~\ref{cp:marco}) frente a los de este trabajo. «---»: sin cifra equivalente.}
    \label{tab:res_estado_arte}
    \footnotesize
    \setlength{\tabcolsep}{3pt}
    \begin{tabularx}{\textwidth}{@{}>{\raggedright\arraybackslash}p{2.6cm}
        >{\raggedright\arraybackslash\hspace{0pt}}X
        >{\raggedright\arraybackslash\hspace{0pt}}p{2.5cm}
        >{\centering\arraybackslash}p{1.9cm}>{\centering\arraybackslash}p{1.9cm}@{}}
        \toprule
        \textbf{Trabajo} & \textbf{Datos y protocolo} & \textbf{Métrica} & \textbf{Trabajo citado}
            & \textbf{Este trabajo} \\
        \midrule
        \multicolumn{5}{@{}l}{\textit{Clasificación en siete clases (Swin-Tiny)}} \\
        \textcite{aboulmira2024skinhealthmate} & \emph{HAM10000}; partición no informada & F1 macro & 0.81 & 0.733 \\
        \textcite{kim2025skinsavvy2} & \emph{HAM10000} ampliado; partición al azar por imagen & Sensibilidad macro
            & 0.780 & 0.692 \\
        \textcite{gessert2020skin} & \emph{ISIC 2019}, otras clases; prueba del desafío; validación por lesión & Sensibilidad macro & 0.636
            & 0.692 \\
        \filasep
        \multicolumn{5}{@{}l}{\textit{Detección de melanoma (Swin-Tiny)}} \\
        \textcite{aboulmira2024skinhealthmate} & Clase más probable & Sensibilidad & 0.49 & 0.623 \\
        \textcite{yao2022single} & \emph{ISIC 2018}; modelo elegido sobre la prueba & Sensibilidad & 0.819 & 0.796 \\
        \textcite{marchetti2023prospective} & Lesiones biopsiadas; umbral de alta sensibilidad & AUC & 0.857
            & 0.921 \\
        \filasep
        \multicolumn{5}{@{}l}{\textit{\glspl{vlm} sin ajuste}} \\
        \textcite{yan2025derm1miccv} & DermLIP en \emph{HAM10000} & Exactitud & 0.682 & 0.716 \\
        \textcite{yan2026dermfmzero} & DermFM-Zero en siete conjuntos; DermLIP obtiene 0.561 & Exactitud balanceada y AUC, en media
            & 0.733 & --- \\
        \textcite{zeng2025mmskin} & LLaVA-Med en \emph{HAM10000} & Exactitud & 0.055 & 0.037--0.114 \\
        \textcite{zeng2025mmskin} & SkinVL, especializado, en \emph{HAM10000} & Exactitud & 0.513 & --- \\
        \textcite{jiang2026multimodal} & Generativos abiertos en consultas reales; tres primeros diagnósticos
            & Exactitud & 0.015--0.134 & --- \\
        \bottomrule
    \end{tabularx}
\end{table}

En la clasificación de siete clases, los resultados de Swin-Tiny quedan por debajo de los dos trabajos que usan
\emph{HAM10000}, pero ninguno de ellos divide las imágenes por lesión: uno no informa cómo las divide, y el otro amplía
el conjunto con imágenes externas y lo divide al azar por imagen. La Sección~\ref{sec:res_fuga} mostró que, en esta
prueba, una partición por imagen sobreestima la sensibilidad macro en unas 0.07, una cifra del mismo orden que la
distancia con el segundo trabajo. Por el contrario, el único trabajo que mantiene juntas las imágenes de una lesión
informa una sensibilidad macro algo menor que la de Swin-Tiny, aunque sobre otro conjunto y con otras clases. La
comparación no permite atribuir toda la distancia a la partición, porque los trabajos difieren también en los datos y
en el entrenamiento, pero sí muestra que una parte puede deberse a ella.

En la detección de melanoma, la tabla muestra que la sensibilidad depende tanto de la regla de decisión como del
modelo. Con la clase más probable, Swin-Tiny detecta más melanomas que el ensamble de
\textcite{aboulmira2024skinhealthmate}, que usa la misma regla. El trabajo de \textcite{yao2022single} obtiene una
sensibilidad algo mayor que la de Swin-Tiny con el umbral de validación, pero elige su modelo sobre la misma prueba
en que lo informa. En el otro extremo, \textcite{marchetti2023prospective} operan con una sensibilidad cercana a la
total a cambio de una especificidad baja, sobre lesiones que ya habían sido biopsiadas. Como cada trabajo usa otra
población y otro punto de operación, la comparación no ordena los modelos, pero justifica que este trabajo fije el
umbral en validación y lo informe junto con la sensibilidad.

En los \glspl{vlm} generativos, los resultados coinciden con la literatura. \textcite{zeng2025mmskin} miden LLaVA-Med
sin ajuste en \emph{HAM10000} con una exactitud del 5.5\,\%, dentro del rango que obtiene aquí, y solo su modelo
especializado, SkinVL, supera el 50\,\%. Del mismo modo, \textcite{jiang2026multimodal} informan exactitudes de entre el 1.5\,\% y el 13.4\,\% con los modelos abiertos en
consultas reales, entre ellos LLaVA-Med. Por lo tanto, el colapso de los generativos de propósito general no es propio
de este trabajo. También DermLIP sin ajuste, con las plantillas y los nombres de clase de sus autores, obtiene en
\emph{HAM10000} una exactitud cercana a la que ellos informan, sobre otra partición de prueba. DermFM-Zero, evaluado
por sus autores junto con DermLIP, lo supera a su vez, lo que indica que la especialización en dermatología sigue
mejorando por la vía contrastiva.

En conjunto, la comparación ubica a Swin-Tiny cerca de los trabajos con un protocolo comparable y por debajo de los que
usan una partición por imagen, y confirma que la dificultad de los \glspl{vlm} generativos para clasificar lesiones de
piel aparece también en la literatura.
